\begin{figure}[ht]
\centering
\begin{tikzpicture}[scale=1.0,
  wk/.style   = {circle, draw=engblack, thick, fill=englime!25,
                 inner sep=0pt, minimum size=8mm, font=\small},
  tk/.style   = {circle, draw=engblue, thick, fill=engblue!15,
                 inner sep=0pt, minimum size=8mm, font=\small},
  free/.style = {thin, draw=enggray},
  matched/.style = {very thick, draw=engred},
]
% Left column: workers. Right column: tasks.
\node[wk] (w1) at (0, 3) {$w_1$};
\node[wk] (w2) at (0, 1.5) {$w_2$};
\node[wk] (w3) at (0, 0) {$w_3$};
\node[tk] (t1) at (4, 3) {$t_1$};
\node[tk] (t2) at (4, 1.5) {$t_2$};
\node[tk] (t3) at (4, 0) {$t_3$};

% Edges: w1-t1, w1-t2, w2-t1, w3-t2, w3-t3.
% Greedy matching (dashed-thin unmatched, red-thick matched) BEFORE augment:
%   matched: w1-t1, w3-t2 ; unmatched vertices: w2, t3.
\draw[matched] (w1) -- (t1);
\draw[free]    (w1) -- (t2);
\draw[free]    (w2) -- (t1);
\draw[matched] (w3) -- (t2);
\draw[free]    (w3) -- (t3);

% The augmenting path t3 - w3 - t2 - w1 - t1 - w2 highlighted.
\node[engorange, font=\scriptsize, anchor=west] at (4.4, 1.5)
  {augmenting path};
\node[engorange, font=\scriptsize, anchor=west] at (4.4, 1.05)
  {$t_3\,w_3\,t_2\,w_1\,t_1\,w_2$};

\node[engred, font=\scriptsize, anchor=east] at (-0.4, 3)
  {matched};
\node[enggray, font=\scriptsize, anchor=east] at (-0.4, 1.5)
  {free edge};
\end{tikzpicture}
\caption{A bipartite graph on workers $\{w_1,w_2,w_3\}$ and tasks
$\{t_1,t_2,t_3\}$. The greedy matching $\{w_1t_1,\,w_3t_2\}$ (red)
leaves $w_2$ and $t_3$ unmatched, and no direct edge joins them. The
alternating path $t_3\,w_3\,t_2\,w_1\,t_1\,w_2$ starts and ends at
unmatched vertices and alternates free and matched edges; flipping
its edges yields the perfect matching $\{w_1t_2,\,w_2t_1,\,w_3t_3\}$
that the greedy start missed. Berge's theorem certifies optimality
by the absence of any further augmenting path, and Konig's theorem
equates the size of this maximum matching with the size of a minimum
set of vertices covering every edge.}
\label{fig:vol02:ch17:bipartite-matching}
\end{figure}
